% Общие стили оформления.
% Возможные варианты значений ищите в описании библиотеки beamer
\usetheme{Pittsburgh}
\usecolortheme{whale}

% \usetheme[secheader]{Boadilla}
% \usecolortheme{seahorse}

% Размер полей слайдов
\setbeamersize{text margin left=1cm,%
               text margin right=1cm}

% выключение кнопок навигации
\beamertemplatenavigationsymbolsempty

% Размеры шрифтов
\setbeamerfont{title}{size=\large}
\setbeamerfont{subtitle}{size=\small}
\setbeamerfont{author}{size=\normalsize}
\setbeamerfont{institute}{size=\small}
\setbeamerfont{date}{size=\normalsize}
\setbeamerfont{bibliography item}{size=\small}
\setbeamerfont{bibliography entry author}{size=\small}
\setbeamerfont{bibliography entry title}{size=\small}
\setbeamerfont{bibliography entry location}{size=\small}
\setbeamerfont{bibliography entry note}{size=\small}
% Аналогично можно настроить и другие размеры.
% Названия классов элементов можно найти здесь
% http://www.cpt.univ-mrs.fr/~masson/latex/Beamer-appearance-cheat-sheet.pdf

% Цвет элементов
\setbeamercolor{footline}{fg=blue}
\setbeamercolor{bibliography item}{fg=black}
\setbeamercolor{bibliography entry author}{fg=black}
\setbeamercolor{bibliography entry title}{fg=black}
\setbeamercolor{bibliography entry location}{fg=black}
\setbeamercolor{bibliography entry note}{fg=black}
% Аналогично можно настроить и другие цвета.
% Названия классов элементов можно найти здесь
% http://www.cpt.univ-mrs.fr/~masson/latex/Beamer-appearance-cheat-sheet.pdf

% Нумеровать список статей
% https://tex.stackexchange.com/a/419506/104425
\setbeamertemplate{bibliography item}{\insertbiblabel}
% или убрать номера
% \setbeamertemplate{bibliography item}{}

% Использовать шрифт с засечками для формул
% https://tex.stackexchange.com/a/34267/104425
\usefonttheme[onlymath]{serif}

% https://tex.stackexchange.com/a/291545/104425
\makeatletter
\def\beamer@framenotesbegin{% at beginning of slide
    \usebeamercolor[fg]{normal text}
    \gdef\beamer@noteitems{}%
    \gdef\beamer@notes{}%
}
\makeatother

% footer презентации
\setbeamertemplate{footline}{
    \leavevmode%
    \hbox{%
        \begin{beamercolorbox}[wd=.333333\paperwidth,ht=2.25ex,dp=1ex,center]{}%
            % И. О. Фамилия, Организация кратко
            \thesisAuthorShort, \thesisOrganizationShort
        \end{beamercolorbox}%
        \begin{beamercolorbox}[wd=.333333\paperwidth,ht=2.25ex,dp=1ex,center]{}%
            % Город, 20XX
            \thesisCity, \thesisYear
        \end{beamercolorbox}%
        \begin{beamercolorbox}[wd=.333333\paperwidth,ht=2.25ex,dp=1ex,right]{}%
            Стр. \insertframenumber{} из \inserttotalframenumber \hspace*{2ex}
        \end{beamercolorbox}}%
    \vskip0pt%
}

% вывод на экран заметок к презентации
\ifnumequal{\value{presnotes}}{0}{}{
    \setbeameroption{show notes}
    \ifnumequal{\value{presnotes}}{2}{
        \setbeameroption{show notes on second screen=\presposition}
    }{}
}

% \centerfloat --- команда из класса memoir, используемого диссертацией.
% В beamer она отсутствует, поэтому определяем её здесь, чтобы окружения
% figure, перенесённые из текста диссертации, работали без изменений.
\makeatletter
\providecommand*{\centerfloat}{%
    \parindent \z@
    \leftskip \z@ \@plus 1fil \@minus \textwidth
    \rightskip\leftskip
    \parfillskip \z@skip}
\makeatother

% -----------------------------------------------------------------------------
% Положения, выносимые на защиту. Единственный источник текста в презентации:
% используется и на слайде <<Положения, выносимые на защиту>>, и на слайдах-трекерах.
% Тексты дословно совпадают с common/characteristic.tex --- при правке положений
% в диссертации их нужно синхронизировать здесь.
% -----------------------------------------------------------------------------
\newcommand{\defposA}{Предложенный алгоритм вписывания проекции объекта на виде сверху в ориентированную ограничивающую рамку, опирающийся на модель L-образной формы видимой границы ТС, обеспечивает прирост среднего коэффициента Жаккара оценки положения окружающих ТС относительно сравниваемых общеизвестных алгоритмов, не использующих нейросетевые алгоритмы для вписывания проекции: $0,337$ против $0,328$ у лучшего из них (прирост $2,7$~\%).}
\newcommand{\defposB}{Среди алгоритмов оценки положения объектов на плоскости движения БТС, применимых на бортовом вычислителе БТС в режиме реального времени, предложенный алгоритм комплексирования визуальных и лидарных данных, не использующий приближение поверхности дороги плоскостью, обеспечивает наибольший средний коэффициент Жаккара: прирост $10,0$~\% по сравнению с алгоритмом, использующим только лидарные данные ($0,439$ против $0,399$), и прирост $30,3$~\% по сравнению с алгоритмом, использующим только визуальные данные ($0,439$ против $0,337$).}
\newcommand{\defposC}{Предложенный алгоритм проекции семантической сегментации проезжей части на плоскость движения БТС на основе комплексирования визуальных и лидарных данных в сочетании с базовым лидарным алгоритмом построения карты проходимости пространства повышает её качество по сравнению с картой, строящейся исключительно по лидарным данным: среднеквадратичная ошибка относительно эталонной карты снижается с $0,231$ до $0,218$ (на $5,6$~\%), а среднеквадратичная ошибка стоимости поиска пути (PathFinding Cost Mean Squared Error, PFC-MSE) --- с $1019,3$ до $924,4$ (на $9,3$~\%).}
\newcommand{\defposD}{Предложенные структура данных для представления карты проходимости пространства и алгоритм, допускающий неблокирующее выполнение операций над картой, в одномерном случае уменьшают константу линейной вычислительной сложности цикла публикации карты в $13,89$ раза по сравнению с базовой реализацией, использующей блокировки при выполнении операций.}

% Кегли: активное положение и приглушённые остальные.
\newcommand{\defposactivefont}{\fontsize{8.5}{10.2}\selectfont\color{black}\everymath{\color{black}}}
\newcommand{\defposdimfont}{\fontsize{6}{7.2}\selectfont\color{gray!85}\everymath{\color{gray!85}}}
\newcommand{\defposfont}[2]{\ifnum#1=#2\defposactivefont\else\ifnum#1=0\defposactivefont\else\defposdimfont\fi\fi}

% \defposlist{N} --- все четыре положения полным текстом;
% N-е выводится основным кеглем, остальные --- уменьшенным и приглушённым.
% \defposlist{0} --- все положения основным кеглем (обычный слайд с положениями).
\newcommand{\defposlist}[1]{%
    \begin{enumerate}
        \item[{\defposfont{#1}{1}1.}] {\defposfont{#1}{1}\defposA\par}
        \item[{\defposfont{#1}{2}2.}] {\defposfont{#1}{2}\defposB\par}
        \item[{\defposfont{#1}{3}3.}] {\defposfont{#1}{3}\defposC\par}
        \item[{\defposfont{#1}{4}4.}] {\defposfont{#1}{4}\defposD\par}
    \end{enumerate}}

\newcommand{\illustrationplaceholder}[2][0.9\linewidth]{%
    \begingroup
    \setlength{\fboxrule}{1pt}%
    \setlength{\fboxsep}{6pt}%
    \fcolorbox{red!60!black}{red!4}{%
        \begin{minipage}{#1}%
            \centering
            {\fontsize{7}{8.4}\selectfont\color{red!60!black}\textit{ЗАГЛУШКА: иллюстрацию требуется отрисовать}}\\[0.4em]
            {\fontsize{7}{8.4}\selectfont #2}
        \end{minipage}}%
    \endgroup}
